% GENERATED FILE. Edit canonical lessons, then run npm run generate:books.
% canonical-source-hash: fnv1a64:967e4938d23eee52
% canonical-lessons: PA-C251-doi, PA-C251-chhikka, PA-C251-chhaba, PA-C251-changer, PA-C251-tamba

\chapter{Spoons and Baskets}
\label{ch:pa-a2-spoons-and-baskets}
\hlchaptermodality{251}

\begin{chapteropening}
\textbf{By the end of this chapter:} \emph{I can name kitchen spoons and baskets.}

Complete the last lesson of chapter 251: I can name kitchen spoons and baskets.
\end{chapteropening}

\section[\texorpdfstring{\pa{ਡੋਈ} (ḍoī)}{ਡੋਈ (ḍoī)}]{\pa{ਡੋਈ} (ḍoī) — a wooden spoon}

\label{lesson:PA-C251-doi}



\begin{quote}
Before the new one: say the Punjabi for a water pitcher, then the Punjabi for brass.
\end{quote}

\subsection*{You'll want to know: \pa{ਡੋਈ}}
\textbf{\pa{ਡੋਈ}} — \emph{ḍoī} — \textquotedblleft{}a wooden spoon\textquotedblright{}.

\begin{quote}
\pa{ਡੋਈ} \pa{ਲੱਕੜ} \pa{ਦੀ} \pa{ਹੈ।}
\end{quote}

\begin{quote}
\emph{ḍoī lakkaṛ dī hai.}
\end{quote}

\begin{quote}
The spoon is made of wood.
\end{quote}

A feminine noun.

\begin{cousinweb}
\textbf{\pa{ਡੋਈ}} (\emph{ḍoī}) is inherited from a reconstructed \emph{*ḍōva}, \textquotedblleft{}spoon\textquotedblright{}; CDIAL lists Punjabi \emph{ḍoī}, \textquotedblleft{}spoon\textquotedblright{} (CDIAL 5573).
\end{cousinweb}

\begin{grammarlens}[title={What you've built}]
One more word for this chapter.
\end{grammarlens}

\subsection*{Your turn}
\begin{itemize}
\raggedright
  \item \emph{Say it:} \emph{ḍoī}
  \item \emph{Say it:} \emph{ḍoī}, once more
  \item \emph{Say it:} \emph{ḍoī lakkaṛ dī hai.} — The spoon is made of wood.
  \item \emph{Recall:} say the Punjabi for a water pitcher, then the Punjabi for brass, then say \emph{ḍoī} again
\end{itemize}

\begin{tcolorbox}[breakable,colback=teal!4,colframe=teal!35!black,title={Before you move on}]
What does \pa{ਡੋਈ} mean? (\textquotedblleft{}A wooden spoon\textquotedblright{}.) Say it once more.
\end{tcolorbox}
\section[\texorpdfstring{\pa{ਛਿੱਕਾ} (chhikkā)}{ਛਿੱਕਾ (chhikkā)}]{\pa{ਛਿੱਕਾ} (chhikkā) — a hanging net (for pots)}

\label{lesson:PA-C251-chhikka}



\begin{quote}
Before the new one: say the Punjabi for brass, then the Punjabi for a wooden spoon.
\end{quote}

\subsection*{You'll want to know: \pa{ਛਿੱਕਾ}}
\textbf{\pa{ਛਿੱਕਾ}} — \emph{chhikkā} — \textquotedblleft{}a hanging net (for pots)\textquotedblright{}.

\begin{quote}
\pa{ਦੁੱਧ} \pa{ਦੀ} \pa{ਹਾਂਡੀ} \pa{ਛਿੱਕੇ} \pa{ਵਿੱਚ} \pa{ਹੈ।}
\end{quote}

\begin{quote}
\emph{dudh dī hāṇḍī chhikke vich hai.}
\end{quote}

\begin{quote}
The milk pot is in the hanging net.
\end{quote}

A masculine noun.

\begin{cousinweb}
\textbf{\pa{ਛਿੱਕਾ}} (\emph{chhikkā}) is inherited from Sanskrit \emph{śikya}, \textquotedblleft{}rope-sling for carrying things\textquotedblright{}; CDIAL lists Punjabi \emph{chikkā}, \textquotedblleft{}a hanging net\textquotedblright{} (CDIAL 12427).
\end{cousinweb}

\begin{grammarlens}[title={What you've built}]
One more word for this chapter.
\end{grammarlens}

\subsection*{Your turn}
\begin{itemize}
\raggedright
  \item \emph{Say it:} \emph{chhikkā}
  \item \emph{Say it:} \emph{chhikkā}, once more
  \item \emph{Say it:} \emph{dudh dī hāṇḍī chhikke vich hai.} — The milk pot is in the hanging net.
  \item \emph{Recall:} say the Punjabi for brass, then the Punjabi for a wooden spoon, then say \emph{chhikkā} again
  \item \emph{Write it:} copy \textbf{\pa{ਛਿੱਕਾ}} once, checking the whole word
\end{itemize}

\begin{tcolorbox}[breakable,colback=teal!4,colframe=teal!35!black,title={Before you move on}]
What does \pa{ਛਿੱਕਾ} mean? (\textquotedblleft{}A hanging net (for pots)\textquotedblright{}.) Say it once more.
\end{tcolorbox}
\section[\texorpdfstring{\pa{ਛਾਬਾ} (chhābā)}{ਛਾਬਾ (chhābā)}]{\pa{ਛਾਬਾ} (chhābā) — a flat basket}

\label{lesson:PA-C251-chhaba}



\begin{quote}
Before the new one: say the Punjabi for a wooden spoon, then the Punjabi for a hanging net.
\end{quote}

\subsection*{You'll want to know: \pa{ਛਾਬਾ}}
\textbf{\pa{ਛਾਬਾ}} — \emph{chhābā} — \textquotedblleft{}a flat basket\textquotedblright{}.

\begin{quote}
\pa{ਛਾਬੇ} \pa{ਵਿੱਚ} \pa{ਸੇਬ} \pa{ਹਨ।}
\end{quote}

\begin{quote}
\emph{chhābe vich seb han.}
\end{quote}

\begin{quote}
There are apples in the flat basket.
\end{quote}

A masculine noun.

\begin{cousinweb}
\textbf{\pa{ਛਾਬਾ}} (\emph{chhābā}) is inherited from Sanskrit \emph{chadman}, \textquotedblleft{}roof\textquotedblright{}; CDIAL lists Punjabi \emph{chābbā}, \textquotedblleft{}flat basket\textquotedblright{} (CDIAL 4981).
\end{cousinweb}

\begin{grammarlens}[title={What you've built}]
One more word for this chapter.
\end{grammarlens}

\subsection*{Your turn}
\begin{itemize}
\raggedright
  \item \emph{Say it:} \emph{chhābā}
  \item \emph{Say it:} \emph{chhābā}, once more
  \item \emph{Say it:} \emph{chhābe vich seb han.} — There are apples in the flat basket.
  \item \emph{Recall:} say the Punjabi for a wooden spoon, then the Punjabi for a hanging net, then say \emph{chhābā} again
  \item \emph{Write it:} copy \textbf{\pa{ਛਾਬਾ}} once, checking the whole word
\end{itemize}

\begin{tcolorbox}[breakable,colback=teal!4,colframe=teal!35!black,title={Before you move on}]
What does \pa{ਛਾਬਾ} mean? (\textquotedblleft{}A flat basket\textquotedblright{}.) Say it once more.
\end{tcolorbox}
\section[\texorpdfstring{\pa{ਚੰਗੇਰ} (chaṅger)}{ਚੰਗੇਰ (chaṅger)}]{\pa{ਚੰਗੇਰ} (chaṅger) — a bread basket}

\label{lesson:PA-C251-changer}



\begin{quote}
Before the new one: say the Punjabi for a hanging net, then the Punjabi for a flat basket.
\end{quote}

\subsection*{You'll want to know: \pa{ਚੰਗੇਰ}}
\textbf{\pa{ਚੰਗੇਰ}} — \emph{chaṅger} — \textquotedblleft{}a bread basket\textquotedblright{}.

\begin{quote}
\pa{ਚੰਗੇਰ} \pa{ਵਿੱਚ} \pa{ਗਰਮ} \pa{ਰੋਟੀਆਂ} \pa{ਹਨ।}
\end{quote}

\begin{quote}
\emph{chaṅger vich garam roṭīāṁ han.}
\end{quote}

\begin{quote}
There are hot rotis in the bread basket.
\end{quote}

A feminine noun.

\begin{cousinweb}
\textbf{\pa{ਚੰਗੇਰ}} (\emph{chaṅger}) is inherited from Sanskrit \emph{caṅgērī}, \textquotedblleft{}basket\textquotedblright{}; CDIAL lists Punjabi \emph{caṅger}, \textquotedblleft{}shallow basket\textquotedblright{} (CDIAL 4565). CDIAL notes that the varied modern forms point to a source outside Indo-Aryan.
\end{cousinweb}

\begin{grammarlens}[title={What you've built}]
One more word for this chapter.
\end{grammarlens}

\subsection*{Your turn}
\begin{itemize}
\raggedright
  \item \emph{Say it:} \emph{chaṅger}
  \item \emph{Say it:} \emph{chaṅger}, once more
  \item \emph{Say it:} \emph{chaṅger vich garam roṭīāṁ han.} — There are hot rotis in the bread basket.
  \item \emph{Recall:} say the Punjabi for a hanging net, then the Punjabi for a flat basket, then say \emph{chaṅger} again
  \item \emph{Write it:} copy \textbf{\pa{ਚੰਗੇਰ}} once, checking the whole word
\end{itemize}

\begin{tcolorbox}[breakable,colback=teal!4,colframe=teal!35!black,title={Before you move on}]
What does \pa{ਚੰਗੇਰ} mean? (\textquotedblleft{}A bread basket\textquotedblright{}.) Say it once more.
\end{tcolorbox}
\section[\texorpdfstring{\pa{ਤਾਂਬਾ} (tāmbā)}{ਤਾਂਬਾ (tāmbā)}]{\pa{ਤਾਂਬਾ} (tāmbā) — copper}

\label{lesson:PA-C251-tamba}



\begin{quote}
Before the new one: say the Punjabi for a flat basket, then the Punjabi for a bread basket.
\end{quote}

\subsection*{You'll want to know: \pa{ਤਾਂਬਾ}}
\textbf{\pa{ਤਾਂਬਾ}} — \emph{tāmbā} — \textquotedblleft{}copper\textquotedblright{}.

\begin{quote}
\pa{ਪੁਰਾਣਾ} \pa{ਲੋਟਾ} \pa{ਤਾਂਬੇ} \pa{ਦਾ} \pa{ਹੈ।}
\end{quote}

\begin{quote}
\emph{purāṇā loṭā tāmbe dā hai.}
\end{quote}

\begin{quote}
The old water pot is made of copper.
\end{quote}

A masculine noun.

\begin{cousinweb}
\textbf{\pa{ਤਾਂਬਾ}} (\emph{tāmbā}) is inherited from Sanskrit \emph{tāmra}, \textquotedblleft{}copper\textquotedblright{}; CDIAL lists Punjabi \emph{tāmbā}, \textquotedblleft{}copper\textquotedblright{} (CDIAL 5779).
\end{cousinweb}

\begin{grammarlens}[title={What you've built}]
One more word for this chapter.
\end{grammarlens}

\subsection*{Your turn}
\begin{itemize}
\raggedright
  \item \emph{Say it:} \emph{tāmbā}
  \item \emph{Say it:} \emph{tāmbā}, once more
  \item \emph{Say it:} \emph{purāṇā loṭā tāmbe dā hai.} — The old water pot is made of copper.
  \item \emph{Recall:} say the Punjabi for a flat basket, then the Punjabi for a bread basket, then say \emph{tāmbā} again
  \item \emph{Write it:} copy \textbf{\pa{ਤਾਂਬਾ}} once, checking the whole word
\end{itemize}

\begin{tcolorbox}[breakable,colback=teal!4,colframe=teal!35!black,title={Before you move on}]
What does \pa{ਤਾਂਬਾ} mean? (\textquotedblleft{}Copper\textquotedblright{}.) Say it once more.
\end{tcolorbox}
